\documentclass[../report.tex]{subfiles}
\begin{document}

\section{Mục tiêu và yêu cầu}

Đề bài yêu cầu thiết kế, mô phỏng và kiểm thử một máy tính 32 bit theo kiến trúc tập lệnh RISC-V RV32I
\cite{riscv-spec}, gồm CPU, bộ nhớ chương trình, bộ nhớ dữ liệu và các khối điều khiển. Phải có hai phiên bản:
một CPU đơn chu kỳ và một CPU pipeline 5 tầng IF, ID, EX, MEM, WB xử lý đúng data hazard và control hazard cơ
bản. Dữ liệu, lệnh, thanh ghi và PC đều rộng 32 bit, có 32 thanh ghi. Chương trình phải chạy từ lệnh nạp vào
bộ nhớ lệnh, không được chỉnh tay giá trị thanh ghi. Công cụ là Logisim-evolution \cite{logisim}.

Nhóm làm ba file mạch, đều mở bằng Logisim-evolution 4.1.0:

\begin{itemize}
  \item \file{circ/RV32I\_SingleCycle.circ}: CPU đơn chu kỳ, trong báo cáo gọi tắt là SC.
  \item \file{circ/RV32I\_Pipeline.circ}: CPU pipeline 5 tầng có forwarding, bypass WB$\to$ID, stall khi
        load-use và flush khi rẽ nhánh, gọi tắt là PL. Đây là phiên bản cuối cùng theo yêu cầu đề.
  \item \file{circ/extra/RV32I\_Pipeline\_NoHazard.circ}: cùng pipeline nhưng tắt forwarding và load-use stall,
        gọi tắt là NH. Bản này chỉ chạy đúng khi assembler chèn lệnh \texttt{nop}; nhóm dùng nó để đo xem phần
        cứng xử lý hazard đem lại bao nhiêu.
\end{itemize}

\section{Đối chiếu yêu cầu với sản phẩm}

Bảng~\ref{tab:req} liệt kê từng yêu cầu của đề và vị trí đáp ứng trong project.

\begin{table}[H]
\centering\small
\caption{Yêu cầu của đề và phần hiện thực tương ứng}\label{tab:req}
\begin{tabularx}{\textwidth}{@{}>{\raggedright\arraybackslash}p{5.2cm}X@{}}
\toprule
Yêu cầu & Hiện thực \\
\midrule
Data, lệnh, thanh ghi, PC 32 bit; 32 thanh ghi & Datapath 32 bit, \sig{RegisterFile} có x0 cố định bằng 0 và x1 đến x31 \\
Lệnh số học, logic (10 lệnh R) & \sig{ALU}, \sig{ALUControl}; chương trình \file{t01\_alu.s} \\
Lệnh immediate (9 lệnh I) & \sig{ImmGen} và \sig{ALU}; \file{t01\_alu.s} \\
LW, SW; tùy chọn LB, LH, LBU, LHU, SB, SH & \sig{LoadUnit}, \sig{StoreUnit}, RAM \sig{DMEM}; \file{t05\_mem.s}. Làm cả phần tùy chọn \\
BEQ, BNE, JAL, JALR; tùy chọn BLT, BGE, BLTU, BGEU & \sig{BranchUnit} và MUX chọn PC kế tiếp; \file{t03\_branch.s}, \file{t04\_jump.s}. Làm cả phần tùy chọn \\
Thêm ngoài danh sách tối thiểu & LUI, AUIPC; ECALL và EBREAK dùng làm lệnh dừng (halt) \\
Phiên bản đơn chu kỳ & \file{circ/RV32I\_SingleCycle.circ} (chương~\ref{ch:single}) \\
Pipeline 5 tầng, xử lý hazard & \file{circ/RV32I\_Pipeline.circ}: \sig{ForwardUnit}, \sig{IDBypass}, \sig{HazardUnit} (chương~\ref{ch:pipe}) \\
Chạy từ lệnh nạp vào bộ nhớ lệnh & ROM \sig{IMEM}; thanh ghi chỉ đổi khi lệnh thực thi \\
Công cụ assembly sang mã máy nạp ROM & \file{tools/rvasm.py}: sinh file hex hoặc nạp thẳng vào ROM trong file \file{.circ} \\
Nâng cao: GPIO/UART & \sig{IOController}: UART TX ra TTY, UART RX từ Keyboard, GPIO ra LED 7 đoạn và 8 LED, GPIO vào 8 công tắc DIP \\
Nâng cao: tăng hiệu năng pipeline & Forwarding từ EX/MEM và MEM/WB, bypass WB$\to$ID, tín hiệu \sig{UseRs1}/\sig{UseRs2} để không stall khi lệnh không đọc thanh ghi \\
Bộ test từng module và toàn CPU & 13 file test vector trong \file{tests/vectors/}; 12 chương trình trong \file{programs/} và 25 chương trình ngẫu nhiên trong \file{programs/fuzz/}; \file{tools/run\_tests.py} (chương~\ref{ch:test}) \\
README & \file{README.md}; phụ lục~\ref{app:guide} \\
\bottomrule
\end{tabularx}
\end{table}

Các mục nâng cao chưa làm là dự đoán nhánh động, cache, nhân chia, ngắt và số thực dấu phẩy động. Chương~\ref{ch:ext}
ghi lại phần mở rộng đã làm, các hạn chế và hướng làm tiếp.

Mã nguồn chia thành \file{circ/} (file mạch), \file{generator/} (Python sinh lại file mạch), \file{programs/}
(chương trình test), \file{tests/vectors/} (test vector từng module), \file{tools/} (assembler, golden model, bộ
chạy test) và \file{latex/} (báo cáo này). Cách mở project và chạy test ở phụ lục~\ref{app:guide}.

\end{document}
